6361: \begin{enonce*}[remark]{Remarque 9.13.2}\etiquette[9.13.2]{I.rem9.13.2}
6362: Un\pageoriginale léger effort supplémentaire doit permettre de
6363: remplacer dans \hyperref[I.prop9.13]{9.13} l'hypothèse que $E$ est stable par
6364: petites limites inductives filtrantes par l'hypothèse que $E$
6365: satisfait à $L$ (\hyperref[I.def9.1]{9.1}), si on suppose $E$ stable par
6366: petites sommes, ou que dans $E$ toute \emph{famille} épimorphique
6367: est épimorphique universelle. Il faut alors, dans la démonstration
6368: de c), choisir $\Pi_\circ$ tel que $E$ satisfasse à $L_{\Pi_\circ}$,
6369: et se borner aux sous-catégories pleines $i$ de $C_{/X}$ qui sont
6370: non seulement filtrantes, mais telles que l'ensemble préordonné $\ob
6371: i$ soit grand devant $\Pi_\circ$. Utilisant 8. et l'hypothèse que
6372: $E$ satisfait à $L_{\Pi_\circ}$, on trouve alors que les $X_i$
6373: existent, et on devrait conclure par une variante convenable de
6374: \hyperref[I.lem9.13.1]{9.13.1}, que le rédacteur n'a pas vérifiée.
6375: \end{enonce*}
6376: 
6377: \begin{enonce*}{Proposition 9.14}\etiquette[9.14]{I.prop9.14}
6378: Soient $f:E\rightarrow F$ un foncteur accessible (\hyperref[I.def9.2]{9.2}) entre deux
6379: catégories munies de filtrations cardinales
6380: $(\Filt^{\Pi}(E))_{\Pi\geq\Pi_\circ}$ et
6381: $(\Filt^{\Pi}(F))_{\Pi\geq\Pi'_\circ}$. Alors il existe un cardinal
6382: $\Pi_1\geq \Sup(\Pi_\circ, \Pi'_\circ)$ tel que pour tout cardinal
6383: $\Pi\geq \Pi_1$, on ait
6384: \begin{equation}
6385: f(\Filt^{\Pi}(E))\subset \Filt^{\Pi^{\Pi'}}(F){\quoi}.\tag{9.14.1}
6386: \end{equation}
6387: En particulier, si l'on a $\Pi=2^c$ avec $c\geq \Pi_1$, d'où
6388: $\Pi^{\Pi_1}=2^{c\Pi_1}=2^c=\Pi$, on a
6389: \begin{equation}
6390: f(\Filt^{\Pi}(E))\subset \Filt^{\Pi}(F){\quoi}.\tag{9.14.2}
6391: \end{equation}
6392: \end{enonce*}
6393: 
6394: Signalons tout de suite le
6395: 
6396: \begin{enonce*}{Corollaire 9.15}\etiquette[9.15]{I.cor9.15}
6397: Soit\pageoriginale $E$ une catégorie, munie de deux filtrations
6398: cardinales\linebreak $(\Filt^{\Pi}(E))_{\Pi\geq\Pi_\circ}$ et
6399: $({\Filt'}^{\Pi}(E))_{\Pi\geq %%
6400: \Pi'_\circ}$. Alors il existe un cardinal $\Pi_1\geq \Sup(\Pi_\circ,
6401: \Pi'_\circ)$ tel que, pour tout cardinal $c\geq\Pi_1$, posant
6402: $\Pi=2^c$, on ait $\Filt^{\Pi}(E)={\Filt'}^{\Pi}(E)$.
6403: \end{enonce*}
6404: 
6405: \noindent \emph{Preuve de} 9.14. Soit $c$ un cardinal tel que
6406: $c\geq\Sup(\Pi_\circ, \Pi'_\circ)$, et tel que $f$ soit
6407: $c$-admissible. Soit d'autre part $\Pi_1\geq c$ tel que l'on ait
6408: \begin{equation}
6409: f(\Filt^c(E))\subset \Filt^{\Pi_1}(F){\quoi}.\tag{$\ast$}
6410: \end{equation}
6411: Il existe un tel $\Pi_1$, grâce au fait que $\Filt^c(E)$ est
6412: équivalente à une petite catégorie (9.12 a)), donc $f(\Filt^c(E))$
6413: l'est également, de sorte qu'on peut appliquer 9.12.1 aux objets de
6414: cette dernière pour trouver une $\Filt^{\Pi}(F)$ qui les contient %%
6415: tous (compte tenu que les $\Filt^{\Pi}(F)$ sont des sous-catégories
6416: pleines). Soit donc $\Pi\geq \Pi_1$, et $X\in \ob \Filt^{\Pi}(E)$,
6417: prouvons que $f(X)\in \Filt^{\Pi^{\Pi_1}}(F)$. Écrivons en effet
6418: $$
6419: X=\varinjlim_I X_i{\quoi},
6420: $$
6421: avec les $X_i\in \ob \Filt^c(E)$, $I$ grand devant $c$, $\card I\leq
6422: \Pi^c\leq \Pi^{\Pi^1}$ (\hyperref[I.def9.12]{9.12}c)). Comme $f$ est $c$-admissible, $I$
6423: grand devant $c$, on a
6424: $$
6425: f(X)\simeq \varinjlim_I f(X_i){\quoi},
6426: $$
6427: et comme $\card I\leq \Pi^{\Pi_1}$ et $f(X_i)\in \ob
6428: \Filt^{\Pi_1}(F)\subset \ob \Filt^{\Pi^{\Pi_1}}(F)$ par $(\ast)$, on
6429: a $f(X)\in \Filt^{\Pi^{\Pi_1}}(F)$ par \hyperref[I.def9.12]{9.12} b), \cqfd
6430: 
6431: \setcounter{subsection}{15} \setcounter{subsubsection}{0}
6432: \subsubsection{}\etiquette[9.15.1]{I.sec9.15.1}
6433: La notion de filtration cardinale \hyperref[I.def9.12]{9.12} n'a guère
6434: d'intérêt que lorsque les objets de $E$ sont accessibles. Signalons
6435: qu'il résulte de \hyperref[I.prop9.11]{9.11} que cette condition est satisfaite
6436: lorsque, en plus des hypothèses de\pageoriginale \hyperref[I.prop9.13]{9.13},
6437: on suppose que le foncteur $\Ker$ sur la catégorie des doubles
6438: flèches de $E$ est accessible. Signalons d'autre part:
6439: 
6440: \begin{enonce*}{Proposition 9.16}\etiquette[9.16]{I.prop9.16}
6441: Soit $E$ une catégorie munie d'une filtration cardinale
6442: $(\Filt^{\Pi}(E))_{\Pi\geq\Pi_\circ}$. Supposons que les éléments de
6443: $\Filt^{\Pi_\circ}(E)$ soient des objets accessibles de $E$; alors
6444: il existe un cardinal $\Pi_1\geq\Pi_\circ$ dans $𝒰$ tel que
6445: les objets de $\Filt^{\Pi_\circ}(E)$ soient $\Pi_1$-accessibles; si
6446: $\Pi_1$ est choisi ainsi, alors pour tout cardinal $\Pi\geq \Pi_1$,
6447: on a (avec les notations de \hyperref[I.sec9.3.1]{9.3.1}):
6448: \begin{equation}
6449: E_{\Pi}\subset \Filt^{\Pi}(E)\subset
6450: E_{(\Pi^{\Pi_\circ})}{\quoi}.\tag{9.16.1}
6451: \end{equation}
6452: \end{enonce*}
6453: 
6454: L'existence de $\Pi_1$ résulte immédiatement du fait que
6455: $\Filt^{\Pi_\circ}(E)$ est équivalente à une petite
6456: catégorie (\hyperref[I.def9.12]{9.12}a)). Soit alors $\Pi\geq \Pi_1$. Si $X$
6457:   est dans $\Filt^{\Pi}(E)$,
6458: écrivant $X\simeq \varinjlim_I X_i$ avec $\card I\leq
6459: \Pi^{\Pi_\circ}$, $X_i\in \ob \Filt^{\Pi_\circ}(E)$ pour tout $i$
6460: (\hyperref[I.def9.12]{9.12}c)), alors il résulte de \hyperref[I.cor9.9]{9.9} que $X$ est
6461: $\Pi^{\Pi_\circ}$-accessible, d'où la deuxième inclusion (9.16.1).
6462: Supposons que $X$ soit $\Pi$-accessible, et écrivons $X\simeq
6463: \varinjlim_I X_i$, avec $I$ grand devant $\Pi$ et les $X_i\in \ob
6464: \Filt^{\Pi}(E)$ (\hyperref[I.def9.12]{9.12}c)). Par hypothèse sur $X$, l'isomorphisme
6465: donné $X\xrightarrow{\sim} \varinjlim_I X_i$ se factorise par un des
6466: $X_i$, donc $X$ est isomorphe à un facteur direct de cet $X_i$. On
6467: en conclut que $X\in \ob \Filt^{\Pi}(E)$, donc la première inclusion
6468: (9.16.1), grâce au
6469: 
6470: \begin{enonce*}{Lemme 9.16.2}\etiquette[9.16.2]{I.lem9.16.2}
6471: Tout objet de $E$ qui est un facteur direct d'un objet de
6472: $\Filt^{\Pi}(E)$ est dans $\Filt^{\Pi}(E)$.
6473: \end{enonce*}
6474: 
6475: En effet, si $X$ est l'image d'un projecteur $p$ dans l'objet $Y$ de
6476: $E$ (\ie d'un endomorphisme $p$ tel que $p^2=p$), et si $I$ est un
6477: ensemble ordonné\pageoriginale  filtrant, $X$ est limite inductive
6478: du système inductif filtrant $(Y_i)_{i\in I}$ défini par $Y_i=Y$
6479: pour tout $i\in I$, $p:Y_i\rightarrow Y_j$ si $i<j$. Prenant $I$
6480: grand devant $\Pi_\circ$ et $\card I\leq \Pi$, on voit donc que si
6481: $Y\in \Filt^{\Pi}(E)$, il en est de même de $X$ en vertu de
6482: \hyperref[I.def9.12]{9.12} b).
6483: 
6484: \begin{enonce*}{Corollaire 9.17}\etiquette[9.17]{I.cor9.17}
6485: Sous les conditions de \hyperref[I.prop9.16]{9.16}, pour tout cardinal $c\geq
6486: \Pi_1$, posant $\Pi=2^c$, $\Filt^{\Pi}(E)$ est identique à la
6487: sous-catégorie strictement pleine de $E_{\Pi}$ de $E$ formée des
6488: objets $\Pi$-accessibles.
6489: \end{enonce*}
6490: 
6491: \begin{enonce*}{Corollaire 9.18}\etiquette[9.18]{I.cor9.18}
6492: Soient $E$ une catégorie satisfaisant la condition $L_{\Pi}$
6493: (\hyperref[I.def9.1]{9.1}), où $\Pi$ est un cardinal, et $C$ une
6494: sous-catégorie pleine
6495: de $E$. On désigne par $\Ind(C)_{\Pi}$ la sous-catégorie pleine de
6496: la catégorie $\Ind(C)$ des ind-objets de $C$ (\hyperref[I.sec8.2]{8.2})
6497: formée des
6498: ind-objets de la forme $(X_i)_{i\in I}$, où $I$ est un ensemble
6499: ordonné grand devant $\Pi$. Considérons le foncteur canonique
6500: \begin{equation}
6501: (X_i)_{i\in I}\longmapsto \varinjlim_I
6502: X_i:\Ind(C)_{\Pi}\longrightarrow E{\quoi}.\tag{9.18.1}
6503: \end{equation}
6504: \begin{enumerate}
6505: \renewcommand{\theenumi}{\alph{enumi}}
6506: \renewcommand{\labelenumi}{\theenumi)}
6507: \item Pour que ce foncteur soit pleinement fidèle, il faut et il
6508: suffit que tout objet $X$ de $C$ soit un objet $\Pi$-accessible
6509: (\hyperref[I.def9.3]{9.3}) de $E$.
6510: 
6511: \item Plaçons-nous sous les conditions de \hyperref[I.cor9.17]{9.17}, en particulier
6512: $\Pi=2^c$, et prenons $C=\Filt^{\Pi}(E)$. Alors le foncteur (9.18.1)
6513: est une équivalence de catégories.
6514: \end{enumerate}
6515: \end{enonce*}
6516: 
6517: L'assertion a) est une généralisation immédiate de \hyperref[I.prop8.7.5]{8.7.5}
6518: a), et se prouve de la même façon. Alors b) résulte de
6519: \hyperref[I.cor9.17]{9.17} et de la condition \hyperref[I.def9.12]{9.12} c) des
6520: filtrations cardinales, qui implique que le foncteur envisagé est
6521: essentiellement surjectif.
6522: 
6523: \begin{enonce*}{Corollaire 9.19}\etiquette[9.19]{I.cor9.19}
6524: Sous\pageoriginale les conditions de \hyperref[I.cor9.17]{9.17}, soient
6525: $C=\Filt^{\Pi}(E)$, $F$ une $𝒰$-catégorie, et considérons le
6526: foncteur
6527: \begin{equation}
6528: \Hhom(E,F)_{\Pi}\longrightarrow \Hhom(C,F)\tag{9.19.1}
6529: \end{equation}
6530: induit par le foncteur « restriction à $C$ » $f\mapsto f| C$),
6531: où la source de (9.19.1) est la catégorie des foncteurs
6532: $\Pi$-accessibles de $E$ dans $F$ (\hyperref[I.def9.2]{9.2}). Le foncteur
6533: précédent est
6534: pleinement fidèle. Si $F$ satisfait à la condition $L_{\Pi}$
6535: (\hyperref[I.def9.1]{9.1}), alors le foncteur (9.19.1) est une équivalence de
6536: catégories.
6537: \end{enonce*}
6538: 
6539: La première assertion se prouve comme \hyperref[I.prop7.8]{7.8}, en utilisant
6540: \hyperref[I.cor9.18]{9.18} b). La deuxième s'obtient en construisant un
6541: foncteur quasi-inverse de (9.19.1), en associant à tout
6542: foncteur $f:C\rightarrow F$ le foncteur $(X_i)_{i\in I}\rightarrow
6543: \varinjlim_i f(X_i)$ de $\Ind(E_{\Pi})$ dans $F$, et en utilisant
6544: \hyperref[I.cor9.18]{9.18} b) pour en déduire un foncteur $\bar{f}:E\rightarrow
6545: F$. Tout revient à %%
6546: montrer que ce%%
6547: dernier est $\Pi$-accessible. Or cela se prouve comme l'assertion
6548: analogue \hyperref[I.prop8.7.3]{8.7.3}.
6549: 
6550: \begin{enonce*}{Corollaire 9.20}\etiquette[9.20]{I.cor9.20}
6551: Soient $E$ une catégorie admettant une filtration cardinale et telle
6552: que tout objet de $E$ soit accessible (\cf \hyperref[I.prop9.16]{9.16}), $F$
6553: une catégorie. Alors:
6554: \begin{enumerate}
6555: \renewcommand{\theenumi}{\alph{enumi}}
6556: \renewcommand{\labelenumi}{\theenumi)}
6557: \item La catégorie $\Hhom(E,F)_{acc}$ des foncteurs accessibles
6558: de $E$ dans $F$ (\hyperref[I.def9.2]{9.2}) est une
6559: $𝒰$-catégorie. ({\rm NB} on
6560: rappelle (\hyperref[I.sec9.0]{9.0}) que les catégories données $E$, $F$
6561: sont supposées
6562: être des $𝒰$-catégories.)
6563: 
6564: \item Supposons que $F$ soit stable par petites limites inductives
6565: filtrantes. Pour toute sous-catégorie pleine $C$ de $E$ équivalente
6566: à une petite catégorie, à foncteur d'inclusion $i:C\rightarrow E$,
6567: considérons le foncteur correspondant
6568: $$
6569: i_{!}:\Hhom(C,F)\longrightarrow \Hhom(E,F)
6570: $$
6571: $(5.1)$.\pageoriginale Pour qu'un foncteur $f:E\rightarrow F$ soit
6572: accessible, il faut et il suffit qu'il existe une petite
6573: sous-catégorie $C$ de $E$, telle que $f$ soit dans l'image
6574: essentielle du foncteur précédent $i_{!}$.
6575: \end{enumerate}
6576: \end{enonce*}
6577: 
6578: \noindent \emph{Démonstration}.
6579: \begin{enumerate}
6580: \renewcommand{\theenumi}{\alph{enumi}}
6581: \renewcommand{\labelenumi}{\theenumi)}
6582: \item Il suffit de prouver que pour tout cardinal $\Pi$ tel que $E$
6583: satisfasse $L_{\Pi}$, la sous-catégorie pleine $\Hhom(E,F)_{\Pi}$ de
6584: $\Hhom(E,F)_{acc}$ est une $𝒰$-catégorie. Il suffit
6585: évidemment de le vérifier pour les cardinaux de la forme $2^c$, avec
6586: $c$ assez grand. Mais alors cela résulte de \hyperref[I.cor9.19]{9.19}, puisque
6587: ($C=\Filt^{\Pi}(E)$ étant essentiellement petite) $\Hhom(C,F)$ est
6588: évidemment une $𝒰$-catégorie.
6589: 
6590: \item Par transitivité de la formation des foncteurs $i_{!}$, on
6591: peut dans l'énoncé se borner aux sous-catégories $C$ de la forme
6592: $\Filt^{\Pi}(E)$, où $\Pi$ est comme dans \hyperref[I.cor9.17]{9.17}. On voit alors
6593: aisément que le foncteur composé $\Hhom(\Filt^{\Pi}(E),F)\rightarrow
6594: \Hhom(E,F)_{\Pi}\rightarrow \Hhom(E,F)$, où la première flèche est
6595: quasi-inverse de (9.19.1), et la deuxième est l'inclusion, n'est
6596: autre que le foncteur $i_{!}$, à isomorphisme près. Donc l'assertion
6597: b) résulte de \hyperref[I.cor9.19]{9.19}.
6598: \end{enumerate}
6599: 
6600: \begin{enonce*}[remark]{*Exercice 9.20.1}\etiquette[9.20.1]{I.exer9.20.1}
6601: (Le présent exercice utilise les notions de site et de topos,
6602: développés dans les exposés II et IV.) Soient $E$ un
6603: $𝒰$-topos, $𝒱$ un univers tel que $𝒰\in
6604: 𝒱$, $C$ une petite sous-catégorie pleine génératrice du
6605: $𝒰$-topos $E$, $\Pi_\circ\in 𝒰$ un cardinal infini
6606: tel que $\Pi_\circ\in \card F\ell\,  C$. Pour tout cardinal $\Pi\geq
6607: \Pi_\circ$, $\Pi\in 𝒰$, soit $\Filt^{\Pi}(E)$ la
6608: sous-catégorie strictement pleine de $E$ formée des objets $X$ tels
6609: qu'il existe une famille épimorphique stricte\pageoriginale
6610: $(X_i\rightarrow X)_{i\in I}$ de but $X$, telle que $\card I\leq
6611: \Pi$ et que $X_i\in \Ob C$ pour tout $i\in I$.
6612: \begin{enumerate}
6613: \renewcommand{\theenumi}{\alph{enumi}}
6614: \renewcommand{\labelenumi}{\theenumi)}
6615: \item Montrer que $(\Filt^{\Pi}(E))_{\Pi\geq \Pi_\circ}$ est une
6616: filtration cardinale de la $𝒰$-catégorie $E$, et qu'on peut
6617: choisir $\Pi_\circ$ tel que pour tout $\Pi\geq \Pi_\circ$, on ait
6618: les inclusions
6619: $$
6620: E_{\Pi}\subset \Filt^{\Pi}(E)\subset E_{\Pi^{\Pi_\circ}}{\quoi},
6621: $$
6622: où pour tout cardinal $c$, $E_c$ désigne la sous-catégorie
6623: strictement pleine de $E$ formée des objets $c$-accessibles.
6624: 
6625: \item Choisissant un cardinal $\Pi\geq \Pi_\circ$ tel que
6626: $\Pi=\Pi^{\Pi_\circ}$ (par exemple $\Pi$ de la forme $2^c$, avec
6627: $c>\Pi_\circ$), et posant $C=\Filt^{\Pi}(E)$, montrer qu'on a une
6628: équivalence de catégories
6629: $$
6630: \Ind(C)_{\Pi}\xrightarrow{\approx} E{\quoi},
6631: $$
6632: (notation $\Ind(C)_{\Pi}$ de \hyperref[I.cor9.18]{9.18}).
6633: 
6634: \item Gardons les notations de b), et soient $𝒱$ un univers tel
6635: que $𝒰\subset 𝒱$, $%%
6636: E^{\sim{\phantom{,}}}_{\underline{V}}$ le $𝒱$-topos de
6637: $𝒱$-faisceaux sur le site $E$ (muni de sa topologie
6638: canonique). Désignons par $\Ind(C,𝒱)_{\Pi}$ la catégorie des
6639: $𝒱\text{-}\Ind$-objets de $C$ indexés par des ensembles%%
6640:  d'indices préordonnés \emph{grands devant} $\Pi$. Montrer qu'on a
6641: une équivalence de catégories
6642: $$
6643: \Ind(C,𝒱)_{\Pi}\xrightarrow{\approx}%%
6644: E^{\tilde{\phantom{,}}}_{𝒱}{\quoi}.
6645: $$
6646: 
6647: \item Soient $c\in 𝒱$ un cardinal, $\Ind(E,𝒱)'_c$ la
6648: sous-catégorie pleine de $\Ind(E,𝒱)$ formée des
6649: $𝒱$-ind-objets de $E$ indexés par un ensemble préordonné qui
6650: est grand devant tout cardinal $<c$. Prenant $c=\card 𝒰$,
6651: montrer\pageoriginale qu'on a une équivalence de catégories
6652: $$
6653: \displaylines{\hfill \Ind(E,𝒱)'_{c}\xrightarrow{\approx}%%
6654: E^{\tilde{\phantom{,}}}_{𝒱}{\quoi}.\hfill{\llap{\textbf{*}}}}
6655: $$
6656: \end{enumerate}
6657: \end{enonce*}
6658: 
6659: \setcounter{subsection}{20}
6660: \subsection{}\etiquette[9.21]{I.sec9.21}
6661: La présente section \hyperref[I.sec9.21]{9.21} développe des préliminaires
6662: techniques pour
6663: la démonstration du théorème \hyperref[I.them9.22]{9.22} ci-dessous,
6664: qui constitue le résultat principal du présent paragraphe
6665: \hyperref[I.sec9]{9}. Soit
6666: \begin{equation}
6667: p:E\rightarrow B\tag{9.21.1}
6668: \end{equation}
6669: un foncteur fibrant, où $B$ est une petite catégorie, et où les
6670: foncteurs images inverses
6671: $$
6672: f^\ast:E_\beta\longrightarrow E_\alpha
6673: $$
6674: associés aux flèches $f:\alpha\rightarrow \beta$ de $B$ sont
6675: accessibles (\hyperref[I.def9.2]{9.2}). En particulier, les catégories
6676: fibres $E_\alpha$
6677: $(\alpha \in \ob B)$ satisfont à la condition $L$
6678: (\hyperref[I.def9.1]{9.1}), donc, $B$
6679: étant petite, il existe un cardinal $\Pi'_\circ\in 𝒰$ tel
6680: que toutes les catégories $E_\alpha$ satisfont à la condition
6681: $L_{\Pi'_\circ}$. Soit $\Pi_\circ\in 𝒰$ un cardinal infini
6682: $>\Pi'_\circ$, de sorte que l'on a
6683: \begin{equation}
6684: \Pi_{\circ}>\Pi'_{\circ}, E_\alpha \mbox{ satisfait }
6685: L_{\prod'_{\circ}} \mbox{ pour tout } \alpha\in \ob B.\tag{9.21.2}
6686: \end{equation}
6687: D'ailleurs, les hypothèses faites impliquent aussitôt l'existence
6688: d'un cardinal $c\in 𝒰$ satisfaisant aux conditions
6689: suivantes:
6690: \begin{equation}
6691: \left\{
6692: \begin{tabular}{rp{6cm}}
6693: a) & $c$ est infini,\\
6694: b) & $c\geq \card F\ell\, B$,\\
6695: c) & pour toute flèche $f:\alpha\rightarrow \beta$ dans  $B$, le
6696: foncteur  $f^\ast:E_\beta\rightarrow E_\alpha$
6697:          est $c$-accessible.
6698: \end{tabular}\right.
6699: \tag{9.21.3}
6700: \end{equation}
6701: Supposons\pageoriginale de plus qu'on puisse trouver, pour chaque
6702: $a\in \ob B$, une sous-catégorie pleine $C_\alpha$ de $E_\alpha$,
6703: satisfaisant les conditions suivantes:
6704: \begin{equation}
6705: \left\{
6706: \begin{tabular}{rp{6cm}}
6707: a) & Pour tout $X\in \ob C_\alpha$, $X$ est $c$-accessible dans
6708: $E_\alpha$ (9.3).\\
6709: b) & Pour tout $X\in \ob E_\alpha$, on peut trouver un isomorphisme
6710: $X\simeq \varinjlim_I X_i$ dans $E_\alpha$, où les $X_i$ sont dans
6711: $C_\alpha$ et où $I$ est un ensemble ordonné grand devant $c$.
6712: \end{tabular}\right.
6713: \tag{9.21.4}
6714: \end{equation}
6715: 
6716: Soit $\Pi$ un cardinal tel que l'on ait
6717: \begin{equation}
6718: \Pi\geq \Pi_\circ\quav\mbox{ donc }
6719: \Pi>\Pi'_\circ{\quoi},\tag{9.21.5}
6720: \end{equation}
6721: et pour tout $\alpha\in \ob B$, soit
6722: \begin{equation}
6723: C^{\Pi}_{\alpha}\subset E_{\alpha}\tag{9.21.6}
6724: \end{equation}
6725: formée des objets qui se peuvent représenter sous la forme
6726: $\varinjlim_I X_i$, où $I$ est un ensemble ordonné grand devant
6727: $\Pi'_\circ$, tel que $\card I\leq \Pi$, et où les $X_i$ sont dans
6728: $C_{\alpha}$. Il est clair alors, grâce à \hyperref[I.var7.5.2]{7.5.2}, que
6729: $C^{\Pi}_{\alpha}$ est essentiellement petite, \ie est équivalente à
6730: une petite catégorie.
6731: 
6732: Soit
6733: \begin{equation}
6734: F=\Hhom_B(B,E)\tag{9.21.7}
6735: \end{equation}
6736: la catégorie des sections de $E$ sur $B$, et soit
6737: \begin{equation}
6738: F^{\Pi}\subset F\tag{9.21.8}
6739: \end{equation}
6740: la sous-catégorie strictement pleine de $F$ formée des sections
6741: $X:\alpha\mapsto X(\alpha)$ telles que pour tout $\alpha\in \ob B$
6742: on ait
6743: $$
6744: X(\alpha)\in C^{\Pi}_{\alpha}{\quoi}.
6745: $$
6746: Il\pageoriginale est clair, les $C^{\Pi}_{\alpha}$ étant
6747: essentiellement petites, qu'il en est de même de la catégorie
6748: $F^{\Pi}$. Nous allons montrer que cette catégorie est génératrice,
6749: et plus précisément:
6750: 
6751: \begin{enonce*}{Lemme 9.21.9}\etiquette[9.21.9]{I.lem9.21.9}
6752: Sous les conditions et avec les notations précédentes, on a ce qui
6753: suit:
6754: \begin{enumerate}
6755: \renewcommand{\theenumi}{\roman{enumi}}
6756: \renewcommand{\labelenumi}{(\theenumi)}
6757: \item Tout objet de $F$ est accessible (\hyperref[I.def9.3]{9.3}). Si $d$ est un cardinal
6758: $\geq \Pi'_\circ$, et si $\alpha\mapsto X(\alpha)$ est un élément de
6759: $F$ tel que pour tout $\alpha$, $X(\alpha)$ soit $d$-accessible dans
6760: $E_{\alpha}$, alors $X$ est $d$-accessible dans $F$.
6761: 
6762: \item Supposons $\Pi\leq c$, ou que $\Pi'_\circ$ soit fini (\ie
6763: les $E_{\alpha}$ stables par petites $\varinjlim$ filtrantes). Alors
6764: tout objet $X$ de $F$ est isomorphe à un objet de la forme
6765: $\varinjlim_I X_i$, où les $X_i$ sont dans $F^{\Pi}$ et où $I$ est
6766: grand devant $\Pi$.
6767: \end{enumerate}
6768: \end{enonce*}
6769: 
6770: Pour prouver (i), notons qu'en vertu de (9.21.4) et de
6771: \hyperref[I.cor9.9]{9.9}, tout
6772: objet de $E_{\alpha}$ est accessible. D'autre part, $F$ satisfait à
6773: la condition $L_{\Pi'_\circ}$, en vertu du
6774: 
6775: 
6776: \begin{enonce*}{Lemme 9.21.10}\etiquette[9.21.10]{I.lem9.21.10}
6777: Soit $p:E\rightarrow B$ un foncteur fibrant, $F=\Hhom_B(B,E)$, $I$
6778: une catégorie, $i\mapsto X_i$ un foncteur de $I$ dans $F$. Pour que
6779: $\varinjlim_I X_i$ soit représentable dans $F$, il suffit que pour
6780: tout $\alpha\in \ob B$, $\varinjlim_I X_i(\alpha)$ soit
6781: représentable dans la catégorie fibre $E_{\alpha}$; lorsqu'il en est
6782: ainsi, alors $\varinjlim_I X_i$ « se calcule argument par argument ».
6783: En particulier, si les catégories fibres satisfont à la
6784: condition $L_{\Pi'_\circ}$ ($\Pi'_\circ$ étant un cardinal donné) il
6785: en est de même de $F$.
6786: \end{enonce*}
6787: 
6788: Nous\pageoriginale laissons le détail de la démonstration (facile)
6789: de \hyperref[I.lem9.21.10]{9.21.10} au lecteur, en nous contentant de remarquer qu'il est
6790: commode d'utiliser le résultat suivant, dont la démonstration est
6791: immédiate:
6792: 
6793: \begin{enonce*}{Corollaire 9.21.10.1}\etiquette[9.21.10.1]{I.cor9.21.10.1}
6794: Avec les hypothèses et notations de \hyperref[I.lem9.21.10]{9.21.10} pour
6795: $p:E\rightarrow B$, et pour $I$, pour tout $\alpha\in \ob B$, le
6796: foncteur d'inclusion $E_{\alpha}\rightarrow E$ commute aux limites
6797: inductives de type $I$.
6798: \end{enonce*}
6799: 
6800: Revenant alors aux conditions générales de \hyperref[I.lem9.21.9]{9.21.9}, soit
6801: $X$ un objet de $F$. Comme pour tout $\alpha \in \ob B$, $X(\alpha)$
6802: est accessible dans $E_{\alpha}$, il existe un cardinal $d\in
6803: 𝒰$ tel que pour tout $\alpha$, $X(\alpha)$ soit
6804: $d$-accessible dans $E_{\alpha}$. On peut choisir $d\geq
6805: \Pi'_{\circ}$, de sorte que $F$ satisfait à $L_d$ en vertu de
6806: \hyperref[I.lem9.21.10]{9.21.10}. Utilisant encore \hyperref[I.lem9.21.10]{9.21.10} pour le
6807: calcul des limites inductives $\varinjlim_I Y_i$ dans $F$, avec $I$
6808: grand devant $d$, on constate aussitôt que $X$ est $d$-accessible.
6809: Cela prouve (i).
6810: 
6811: Nous allons prouver maintenant \hyperref[I.lem9.21.9]{9.21.9} (ii) en plusieurs étapes
6812: ((9.21.11) à (\hyperref[I.lem9.21.16]{9.21.16})).
6813: 
6814: Soit donc
6815: $$
6816: X:\alpha\longmapsto X(\alpha)
6817: $$
6818: un objet de $F$. En vertu de (9.21.4 b)), on peut, pour tout
6819: $\alpha\in \ob B$, trouver un ensemble ordonné $I_{\alpha}$ grand
6820: devant $c$, et un isomorphisme $X(\alpha)=\varinjlim_{i\in
6821: I_{\alpha}} X(\alpha)_i$, où $i\mapsto X(\alpha)_i$ est un système
6822: inductif de type $I_{\alpha}$ dans $E_{\alpha}$, les $X(\alpha)_i$
6823: dans $C_{\alpha}$. Considérons l'ensemble ordonné produit
6824: $I=\Pi_{\alpha} I_{\alpha}$; il est clair qu'il est grand devant
6825: $c$, et que les systèmes inductifs\pageoriginale précédents donnent
6826: naissance à des systèmes inductifs dans les $E$,
6827: \begin{equation}
6828: i\longmapsto X(\alpha)_i\in \ob C_{\alpha}, i\in \ob
6829: I{\quoi},\tag{9.21.11}
6830: \end{equation}
6831: indexés par le même ensemble ordonné $I$, et à des isomorphismes
6832: \begin{equation}
6833: X(\alpha)\xrightarrow{\sim}\varinjlim_I X(\alpha)_i\mbox{ dans }
6834: E_{\alpha}{\quoi},\tag{9.21.12}
6835: \end{equation}
6836: pour tout $\alpha\in \ob B$. Nous supposons fixées par la suite des
6837: données (9.21.11) et (9.21.12).
6838: 
6839: \begin{enonce*}{Lemme 9.21.13}\etiquette[9.21.13]{I.lem9.21.13}
6840: Sous les conditions précédentes, on peut trouver une application
6841: $\varphi:I\rightarrow I$ telle que $\varphi(i)\geq i$ pour tout
6842: $i\in I$, et une application $(i,f)\mapsto \lambda(i,f)$ de $I\times
6843: %%
6844: \mathrm{F}{}l\, B$ dans $%%
6845: \mathrm{F}{}l\, E$, satisfaisant aux conditions suivantes:
6846: \begin{enumerate}
6847: \renewcommand{\theenumi}{\alph{enumi}}
6848: \renewcommand{\labelenumi}{\theenumi)}
6849: \item Pour $i\in I$, $(f:\alpha\rightarrow \beta)\in F\ell B$,
6850: $\lambda(i,f)$ est un $f$-morphisme de $X(\alpha)_i$ dans
6851: $X(\beta)_{\varphi(i)}$, rendant commutatif le diagramme
6852: \[
6853: \xymatrix@C=3cm@R=.8cm{
6854:  X(\alpha)\ar[r]^{X(f)} & X(\beta)\\
6855:            &  X(\beta)_{\varphi(i)}\ar[u]\\
6856: X(\alpha)_{i} \ar[uu]\ar[ru]^{\lambda(i,f)}&
6857:                   \\
6858: \alpha \ar[r]^{f}& \beta\quad , }
6859: \]
6860: où les flèches verticales sont les morphismes canoniques déduits de
6861: (9.21.12).
6862: 
6863: \item Pour\pageoriginale tout $i\in I$, et $(f:\alpha\rightarrow
6864:   \beta)\in %%
6865: \mathrm{F}{}l\, B$, on a commutativité dans le diagramme
6866: \[
6867: \xymatrix@C=3cm@R=.9cm@L=6pt{
6868:                       & X(\beta)_{\varphi^{2}(i)}\\
6869:   X(\alpha)_\varphi(i) \ar[ru]^{\lambda(\varphi(i),f)}  & X(\beta)_{\varphi(i)} \ar[u]\\
6870: X(\alpha)_i  \ar[ru]_{\lambda(i,f)} \ar[u]         \\
6871: \alpha \ar[r]^{f} & \beta\quad . }
6872: \]
6873: 
6874: \item Pour tout couple de flèches consécutives
6875: $\alpha\xrightarrow{f}\beta\xrightarrow{g}\gamma$ de $B$, on a
6876: commutativité dans le diagramme suivant
6877: \[
6878: \xymatrix@C=2cm@L=6pt@R=.5cm{
6879: X(\alpha) \ar[r]^{X(f)}& X(\beta) \ar[r]^{X(g)}            & X(\gamma)\\
6880:           &                      & X(\gamma)_{\varphi^{2}(i)}\ar[u]\\
6881:           & X(\beta)_{\varphi(i)} \ar[uu]\ar[ru]^{\lambda(\varphi(i),g)} & X(\gamma)_{\varphi(i)}\ar[u]\\
6882: X(\alpha)_i \ar[uuu]\ar[ru]^{\lambda(i,f)} \ar[rru]_{\lambda(i,gf)}&
6883:  \\
6884: \alpha \ar[r]^{f}& \beta \ar[r]^{g}& \gamma\quad , }
6885: \]
6886: où les flèches verticales sont celles déduites de
6887: (9.21.12).
6888: \end{enumerate}
6889: \end{enonce*}
6890: 
6891: On\pageoriginale notera d'ailleurs que la commutativité des deux
6892: trapèzes supérieurs dans b) est déjà contenu dans a), de sorte que
6893: b) affirme en fait la commutativité du triangle inférieur du
6894: diagramme envisagé.
6895: 
6896: Prouvons \hyperref[I.lem9.21.13]{9.21.13}. Soient $i\in I$, et $f:\alpha \rightarrow \beta$
6897: une flèche dans $B$. Considérons le composé $X(\alpha_i)\rightarrow
6898: X(\alpha)\xrightarrow{X(f)}X(\beta)\simeq \varinjlim_j X(\beta)_j$,
6899: il définit un morphisme dans $E$:
6900: $$
6901: X(\alpha)_i\rightarrow f^\ast (X(\beta))\simeq f^\ast (\varinjlim_j
6902: X(\beta)_j)\simeq \varinjlim_j f^\ast(X(\beta)_j){\quoi},
6903: $$
6904: où la dernière égalité provient du fait que $f^\ast$ est
6905: $c$-accessible (9.21.3 c)%%
6906:  et que $I$ est grand devant $c$. En vertu
6907: de (9.21.4 a)), comme $X(\alpha)\in \ob C_{\alpha}$, le morphisme
6908: envisagé se factorise par un $f^\ast(X(\alpha)_j)$, où $j$ \textit{a
6909: priori} dépend de $i$ et de $f\in %%
6910: \mathrm{F}{}l\,(B)$. Mais en vertu de (9.21.3 b)), et comme $I$ est
6911: grand devant c), on peut choisir $j$ indépendant de $f$, soit
6912: $j=\varphi(i)$. $I$ étant filtrant, on peut supposer $\varphi(i)\geq
6913: i$. On trouve ainsi, pour $i\in I$ et $f\in %%
6914: \mathrm{F}{}l\, F$, un morphisme
6915: $$
6916: X(\alpha)_i\longrightarrow f^\ast(X(\beta)_{\varphi(i)}){\quoi},
6917: $$
6918: ou ce qui revient au même, un $f$-morphisme
6919: $$
6920: \lambda(i,f):X(\alpha)_i\longrightarrow
6921: X(\beta)_{\varphi(i)}{\quoi},
6922: $$
6923: qui par construction rend commutatif le diagramme
6924: \[\vcenter{
6925: \xymatrix@C=3cm{
6926: X(\alpha)\ar[r]^{X(f)} & X(\beta)\\
6927:           & X(\beta)_{\varphi(i)} \ar[u]\\
6928: X(\alpha)_i \ar[uu] \ar[ru]^{\lambda(i,f)}&\\
6929: \alpha \ar[r]^{f} & \beta \quad . }}\leqno{D(f)}
6930: \]
6931: Considérons\pageoriginale alors, pour $i$, $f$, $g$ donnés, le
6932: triangle inférieur du diagramme envisagé dans 9.21.13
6933: c); il n'est
6934: pas clair qu'il est commutatif, mais les deux composés
6935: $X(\alpha)_i\rightrightarrows X(\beta)_{\varphi^2(i)}$ deviennent
6936: égaux après composition avec $X(\beta)_{\varphi^2(i)}\rightarrow
6937: X(\beta)$, comme il résulte de la commutativité des deux trapèzes
6938: supérieurs, et du trapèze contour global, qui sont les diagrammes
6939: $D(f)$, $D(g)$ et $D(gf)$ respectivement. Donc, comme
6940: $X(\alpha)_{\varphi^2(i)}$ est $c$-accessible (9.21.4 a)) et que
6941: $X(\alpha)\simeq \varinjlim_{j\in I}X(\alpha)_j$, avec $I$ grand
6942: devant $c$, il s'ensuit qu'on peut trouver un élément $j\geq
6943: \varphi^2(i)$ de $I$, tel que les deux flèches envisagées deviennent
6944: égales après composition avec $X(\beta)_{\varphi^2(i)}\rightarrow
6945: X(\beta)_j$. A priori, $j$ dépend de $i$, $f$ et $g$. Mais pour $i$
6946: fixé, l'ensemble des couples possibles $f$, $g$ est de cardinal
6947: $\geq c^2=c$, en vertu de (9.21.3 a) et b)), donc, $I$ étant grand
6948: devant $c$, on peut choisir $j$ indépendant de $f$ et de $g$, soit
6949: $j=\varphi'(i)\geq \varphi^2(i)\geq \varphi(i)$. Soit alors, pour
6950: tout $i\in I$ et $f:\alpha\rightarrow \beta$,
6951: $$
6952: \lambda'(i,f):X(\beta)_i\longrightarrow X(\beta)_{\varphi'(i)}
6953: $$
6954: le $f$-morphisme composé $X(\alpha)_i\rightarrow
6955: X(\beta)_{\varphi(i)}\rightarrow X(\beta)_{\varphi'(i)}$, où la
6956: deuxième flèche est le morphisme de transition. Il est alors
6957: immédiat, par construction, que $(\varphi',\lambda')$ satisfait les
6958: conditions 9.21.13 a) et c), pour $(\varphi,\lambda)$. Procédant de
6959: même pour la condition b), on voit qu'on peut choisir $\varphi'$ de
6960: telle façon que cette condition soit également satisfait pour
6961: $(\varphi',\lambda')$. Cela achève la preuve de (9.21.13).
6962: 
6963: \setcounter{subsection}{21} \setcounter{subsubsection}{13}
6964: \subsubsection{}\etiquette[9.21.14]{I.sec9.21.14}
6965: Soit\pageoriginale maintenant
6966: $$
6967: J\subset I
6968: $$
6969: une partie de $I$ satisfaisant les conditions suivantes:
6970: \begin{enumerate}
6971: \renewcommand{\theenumi}{\alph{enumi}}
6972: \renewcommand{\labelenumi}{\theenumi)}
6973: \item $J$ est filtrante,
6974: 
6975: \item pour tout $j\in J$, on a $\varphi(j)\in J$,
6976: 
6977: \item Pour tout $a\in \ob B$, $X_J(\alpha)=\varinjlim_{i\in
6978: J}X(\alpha)_i$ est représentable dans $E_{\alpha}$.
6979: \end{enumerate}
6980: Les conditions a) et c) sont satisfaites en particulier si $J$ est
6981: grand devant $\Pi'_\circ$ (9.21.2). Il résulte alors de
6982: \hyperref[I.cor9.21.10.1]{9.21.10.1}
6983: que pour tout $\alpha\in \ob B$, $X_J(\alpha)$ est la limite
6984: inductive $\varinjlim_{i\in J}X(\alpha)_i$ dans $E$. Or pour une
6985: flèche $f:\alpha\rightarrow \beta$ de $B$, les flèches
6986: $\lambda(i,f)$, pour $i$ variable dans $J$, définissent grâce à
6987: (\hyperref[I.lem9.21.13]{9.21.13} b)) un morphisme de ind-objets de $(X(\alpha)_i)_{i\in J}$
6988: dans $(X(\beta)_i)_{i\in J}$, d'où un homomorphisme
6989: $$
6990: X_J(f):X_J(\alpha)\rightarrow X_J(\beta)
6991: $$
6992: sur les limites inductives, qui est manifestement un
6993: $f$-homomorphisme. Utilisant (\hyperref[I.lem9.21.13]{9.21.13} (c)), on trouve que l'on a des
6994: relations de transitivité
6995: $$
6996: X_J(gf)=X_J(g)X_J(f){\quoi},
6997: $$
6998: de sorte que l'on à défini une section $X_J\in \ob F$ de $E$ sur
6999: $B$. Enfin (\hyperref[I.lem9.21.13]{9.21.13} a)) nous montre que les homomorphismes
7000: canoniques
7001: $$
7002: X_J(\alpha)\longrightarrow X(\alpha)\quav \alpha\in \ob B{\quoi},
7003: $$
7004: sont fonctoriels en $\alpha$, de sorte qu'on a un homomorphisme
7005: canonique
7006: $$
7007: u_J:X_J\longrightarrow X{\quoi}.
7008: $$
7009: D'autre\pageoriginale part, si
7010: $$
7011: J'\supset J
7012: $$
7013: est une autre partie de $I$ satisfaisant aux conditions a), b), c)
7014: ci-dessus, on trouve un homomorphisme canonique
7015: $$
7016: u_{J',J}:X_J\longrightarrow X_{J'}{\quoi},
7017: $$
7018: de sorte que les $X_J$ forment un système inductif dans $F$,
7019: paramétré par l'ensemble (ordonné par inclusion) des parties $J$ de
7020: $I$ satisfaisant les conditions envisagées. Enfin, les
7021: homomorphismes $u_J$ ci-dessus définissent un homomorphismes de ce
7022: système inductif dans (le système inductif constant défini par) $X$.
7023: 
7024: 
7025: \setcounter{subsection}{21} \setcounter{subsubsection}{14}
7026: \subsubsection{}\etiquette[9.21.15]{I.sec9.21.15}
7027: Soit maintenant $K$ un ensemble de parties $J$ de $I$, satisfaisant
7028: aux conditions a), b), c) envisagées dans \hyperref[I.sec9.21.14]{9.21.14}, et
7029: supposons que
7030: $K$ soit filtrant, et de réunion $I$. Alors il est clair que l'on a
7031: $$
7032: \varinjlim_{J\in K}X_J\xrightarrow{\sim}X{\quoi},
7033: $$
7034: en utilisant \hyperref[I.cor9.21.10.1]{9.21.10.1} qui nous ramène à
7035: vérifier qu'on a un
7036: isomorphisme argument par argument.
7037: 
7038: Prenons par exemple pour $K$ l'ensemble de toutes les parties $J$ de
7039: $I$ qui sont grandes devant $\Pi'_\circ$, stables par $\varphi$, et
7040: telles que $\card J\leq \Pi$. Alors par définition (9.21.6) de
7041: $C^{\Pi}_{\alpha}$, on a, pour tout $\alpha\in \ob B$,
7042: $X_J(\alpha)\in \ob C^{\Pi}_{\alpha}$, donc $X_J\in \ob F^{\Pi}$.
7043: Par suite, \hyperref[I.lem9.21.9]{9.21.9} (ii) sera prouvé si nous
7044: établissons que $K$ est
7045: grand devant $\Pi$ (donc filtrant) et de réunion $I$. Il suffira
7046: évidemment, pour ceci, de prouver que toute partie $S$ de $I$ telle
7047: que $\card S\leq \Pi$ est contenue dans une $J\in K$. Ceci
7048: résultera\pageoriginale en effet de l'hypothèse
7049: préliminaire énoncée
7050: dans \hyperref[I.lem9.21.9]{9.21.9} (ii), et du
7051: 
7052: \begin{enonce*}{Lemme 9.21.16}\etiquette[9.21.16]{I.lem9.21.16}
7053: Soient $I$ un ensemble ordonné, $\Pi$ un cardinal infini,
7054: $\varphi:I\rightarrow I$ une application de $I$ dans lui-même.
7055: \begin{enumerate}
7056: \renewcommand{\theenumi}{\alph{enumi}}
7057: \renewcommand{\labelenumi}{\theenumi)}
7058: \item Supposons $I$ filtrant. Alors pour toute partie $S$ de $I$ telle
7059: que $\card S\leq \Pi$, il existe une partie filtrante $J\supset S$
7060: de $I$, telle que $\varphi(J)\subset J$ et $\card(J)\leq \Pi$.
7061: 
7062: \item Soit $\Pi'_{\circ}$ un cardinal $<\Pi$, et supposons $I$ grand
7063: devant $\Pi$. Alors pour toute partie $S$ de $I$ telle que $\card
7064: S\leq \Pi$, il existe une partie $J\supset S$ de $I$ grande devant
7065: $\Pi'_{\circ}$, telle que $\varphi(J)\subset J$ et $\card J\leq
7066: \Pi$.
7067: \end{enumerate}
7068: \end{enonce*}
7069: 
7070: Prouvons par exemple b) (la démonstration de a) étant analogue et
7071: plus simple). Soit $P$ l'ensemble des parties de $I$ de cardinal
7072: $\leq \Pi$, et soit $T\mapsto i_T:P\rightarrow I$ une application
7073: telle que pour $T\in P$, $i_T$ soit un majorant de $T$ dans $I$;
7074: l'existence de cette application exprime simplement l'hypothèse que
7075: $I$ est grand $\Pi$. Soit $A$ un ensemble bien ordonné tel que
7076: $\card A=\Pi$, et tel que pour tout $a\in A$, l'ensemble des $a'\leq
7077: a$ soit de cardinal $<\Pi$. Il s'ensuit en particulier, comme
7078: $\Pi'_{\circ}<\Pi$, que $A$ \emph{est grand devant} $\Pi'_{\circ}$.
7079: Définissons par récurrence transfinie des applications
7080: $$
7081: S\longmapsto S_a: P\longrightarrow P\quad(a\in A){\quoi},
7082: $$
7083: par les formules suivantes:
7084: 
7085: $S_{\circ}=S$, $S_{a+1}=S_a\cup \varphi(S_a)\cup\{i_{S_a}\}$,
7086: $S_a=\bigcup_{a<a}S_a$ si a un ordinal limite.
7087: 
7088: Posons\pageoriginale alors, pour $S\in P$,
7089: $$
7090: J(S)=\bigcup_{a\in A}S_a{\quoi}.
7091: $$
7092: Il est clair alors que $J(S)\supset S$, que $J(S)$ est stable par
7093: $\varphi$, que $\card J(S)\leq \Pi$ (puisque $\card A=\Pi$), enfin
7094: que $J(S)$ est grand devant $\Pi'_{\circ}$ (en utilisant le fait que
7095: $A$ est grand devant $\Pi'_{\circ}$,
7096: 
7097: Cela démontre \hyperref[I.lem9.21.16]{9.21.16} et achève la
7098: démonstration de \hyperref[I.lem9.21.9]{9.21.9}.
7099: 
7100: Signalons aussi une variante de \hyperref[I.lem9.21.9]{9.21.9}:
7101: 
7102: \begin{enonce*}{Lemme 9.21.17}\etiquette[9.21.17]{I.lem9.21.17}
7103: Les notations sont celles de \hyperref[I.lem9.21.9]{9.21.9}. On désigne par $F'$ la
7104: sous-catégorie pleine $\Hhomcart_B(B,E)$ de $F=\Hhom_B(B,E)$ formée
7105: des sections cartésiennes de $E$ sur $F$. Alors:
7106: \begin{enumerate}
7107: \renewcommand{\theenumi}{\roman{enumi}}
7108: \renewcommand{\labelenumi}{(\theenumi)}
7109: \item Tout objet de $F'$ est accessible. Plus précisément, soit
7110: $\alpha\mapsto X(\alpha)$ un élément de $F'$, et soit $d$ un
7111: cardinal tel que $d\geq \Pi'_{\circ}$, $d\geq c$, et tel que pour
7112: tout $\alpha\in \ob B$, $X(\alpha)$ soit un objet $d$-accessible de
7113: $E$. Alors $X$ est un objet $d$-accessible de $F'$.
7114: 
7115: \item Supposons $\Pi\leq c$ et que les foncteurs
7116: $f^\ast:E_{\beta}\rightarrow E_{\alpha}$ soient
7117: $\Pi'_{\circ}$-accessibles, ou que les catégories $E_{\alpha}$
7118: soient stables par petites limites inductives filtrantes, et que les
7119: foncteurs $f^\ast:E_{\beta}\rightarrow E_{\alpha}$ commutent aux
7120: dites limites inductives. Alors tout objet $X$ de $F'$ est isomorphe
7121: à un objet de la forme $\varinjlim_I X_i$, où pour tout $i\in I$,
7122: $X_i$ est un objet de la sous-catégorie pleine ${F'}^{\Pi}=F'\cap
7123: F^{\Pi_{\circ}}$ de $F'$, et où $I$ est grand devant $\Pi$.
7124: \end{enumerate}
7125: \end{enonce*}
7126: 
7127: La\pageoriginale démonstration étant toute analogue à celle de
7128: \hyperref[I.lem9.21.9]{9.21.9}, nous nous contentons d'indiquer les points où une
7129: modification de cette dernière est nécessaire. On note d'abord:
7130: 
7131: \begin{enonce*}{Lemme 9.21.18}\etiquette[9.21.18]{I.lem9.21.18}
7132: L'énoncé \hyperref[I.lem9.21.10]{9.21.10} reste valable lorsqu'on $y$ remplace
7133: $F=\Hhom_B(B,E)$ par $F'=\Hhomcart_B(B,E)$, pourvu que l'on suppose
7134: que les foncteurs images inverse $f^\ast:E_{\beta}\rightarrow
7135: E_{\alpha}$ commutent aux limites inductives de type $I$.
7136: \end{enonce*}
7137: 
7138: Si donc sous les conditions de \hyperref[I.lem9.21.17]{9.21.17}, $d$ est un cardinal tel que
7139: $d\geq c$ et $d\geq \Pi'_{\circ}$, il résulte de ce qui précède et
7140: de l'hypothèse (9.21.3 c)) que pour tout ensemble
7141: ordonné $I$ grand
7142: devant $d$, les limites inductives de type $I$ sont représentables
7143: dans $F'$ et se calculent argument par argument, d'où la conclusion
7144: \hyperref[I.lem9.21.17]{9.21.17} (i) par un argument immédiat.
7145: 
7146: Pour prouver (ii), il faut donner un complément à \hyperref[I.lem9.21.13]{9.21.13}:
7147: 
7148: \begin{enonce*}{Lemme 9.21.19}\etiquette[9.21.19]{I.lem9.21.19}
7149: Sous les conditions de \hyperref[I.lem9.21.13]{9.21.13}, supposons que $X\in \ob F'$ \ie
7150: que $X$ soit une section cartésienne de $E$ sur $B$. Alors on peut
7151: renforcer la conclusion par l'assertion qu'il existe une fonction
7152: $μ$ qui, à tout $i\in I$ et toute flèche $f:\alpha\rightarrow
7153: \beta$ de $B$, associe un morphisme $μ(i,f):f^\ast
7154: X(\beta)_i\rightarrow X(\alpha)_{\varphi(i)}$ dans $E_{\alpha}$, de
7155: telle façon que l'on ait:
7156: 
7157: \begin{enumerate}
7158: \renewcommand{\theenumi}{\alph{enumi}}
7159: \renewcommand{\labelenumi}{\theenumi)}
7160: \setcounter{enumi}{3}
7161: \item Pour tout $i\in I$ et toute flèche $f:\alpha\rightarrow \beta$
7162: de $B$, les deux diagrammes suivants sont commutatifs:
7163: 
7164: \begin{displaymath}
7165: \xymatrix@L=5pt@R=1.25cm@C=2.2cm{
7166:                        & f^{\ast}X(\beta)_{\varphi^{2}(i)}\\
7167: X(\alpha)_{\varphi(i)}\ar[ru]^{\lambda(\varphi (i),f)} &\\
7168:                        & f^{\ast}X(\beta)_i \ar[lu]^{\mu(i,f)}\ar[uu]
7169: } \qquad\qquad \xymatrix@L=5pt@R=1.35cm@C=2cm{
7170: X(\alpha)_{\varphi^2(i)} & \\
7171:                          & f^{\ast}X(\alpha)_{\varphi(i)} \ar[lu]_{\mu(\varphi(i),f)}\\
7172: x(\alpha)_i \ar[uu]\ar[ru]_{\lambda(i,j)}             &. }
7173: \end{displaymath}
7174: \end{enumerate}
7175: \end{enonce*}
7176: 
7177: Pour\pageoriginale le prouver, on procède comme dans 9.21.3, en
7178: considérant le morphisme composé $f^\ast(X(\beta)_i)\rightarrow
7179: f^\ast(X(\beta))\xrightarrow{X(f)^{-1}}X(\alpha)$, où (comme déjà
7180: dans l'écriture de la condition d)) on identifie dans les notations
7181: un $f$-morphisme $R\rightarrow S$ de $E$ avec le morphisme
7182: correspondant $R\rightarrow f^\ast(S)$ de $E_{\alpha}$. Comme
7183: $X(\alpha)=\varinjlim_j X(\alpha)_j$, et $I$ grand devant $c$, on
7184: peut factoriser le morphisme précédent par un des $X(\alpha)_j$, où
7185: $j$ \textit{a priori} dépend de $i$ \emph{et} de $f$, mais peut être
7186: choisi indépendant de $f$, soit $\psi(i)$. Quitte à agrandir la
7187: fonction $\varphi$ de \hyperref[I.lem9.21.13]{9.21.13}, on peut supposer que $\psi=\varphi$.
7188: En procédant comme pour les conditions b) et c) de
7189: \hyperref[I.lem9.21.13]{9.21.13}, on voit
7190: que, quitte à agrandir encore l'application $\varphi$, on peut
7191: supposer que les deux diagrammes de \hyperref[I.lem9.21.19]{9.21.19} d) sont commutatifs.
7192: 
7193: \setcounter{subsection}{21} \setcounter{subsubsection}{19}
7194: \subsubsection{}\etiquette[9.21.20]{I.sec9.21.20}
7195: L'application $\varphi$ étant choisie comme dans
7196: \hyperref[I.lem9.21.19]{9.21.19}, reprenons
7197: l'argument de 9.21.14, où il faut cependant supposer que $J$
7198: satisfait, en plus des conditions énoncées a) b) c), à la condition:
7199: \begin{enumerate}
7200: \renewcommand{\theenumi}{\alph{enumi}}
7201: \renewcommand{\labelenumi}{\theenumi)}
7202: \setcounter{enumi}{3}
7203: \item Pour tout flèche $f:\alpha\rightarrow \beta$ dans $B$, le
7204: foncteur $f^\ast:E_{\beta}\rightarrow E_{\alpha}$ commute aux
7205: limites inductives de type $J$.
7206: \end{enumerate}
7207: 
7208: Je\pageoriginale dis que, moyennant cette condition supplémentaire,
7209: la section $X_J$ de $E$ sur $B$ est cartésienne, \ie pour toute
7210: flèche $f:\alpha\rightarrow \beta$ de $B$, le morphisme
7211: \begin{equation}
7212: X_J(\alpha)\longrightarrow f^\ast X_J(\beta)\tag{$\ast$}
7213: \end{equation}
7214: est un isomorphisme. En effet, en vertu de la condition d)
7215: ci-dessus, le but du morphisme envisagé s'identifie à
7216: $\varinjlim_{i\in J}f^\ast(X(\beta)_i)$, et le morphisme s'obtient
7217: par passage aux $\varinjlim_J$ à partir du morphisme de ind-objets
7218: dans $E_{\alpha}$
7219: $$
7220: (X(\alpha)_i)_{i\in J}\longrightarrow (f^\ast)(X(\beta)_i)_{i\in
7221: J}{\quoi},
7222: $$
7223: déduit des morphismes $\lambda(i,f):X(\alpha)_i\rightarrow
7224: f^\ast(X(\beta)_{\varphi(i)})$ pour $i\in J$. Or les conditions
7225: explicitées dans \hyperref[I.lem9.21.19]{9.21.19} nous assurent que
7226: l'homomorphisme précédent
7227: de ind-objets est en fait un isomorphisme, un inverse étant obtenu
7228: par l'homomorphisme déduit du système des $\mu(i,f)$, pour $i\in J$.
7229: Cela prouve notre assertion que $(\ast)$ est un %%
7230: isomorphisme, donc que $X\in \ob F'$.
7231: 
7232: \subsubsection{}\etiquette[9.21.21]{I.sec9.21.21}
7233: Nous allons supposer maintenant que les foncteurs
7234: $f^\ast:E_{\beta}\rightarrow E_{\alpha}$ sont
7235: $\Pi'_{\circ}$-\emph{accessibles.} Alors les conditions sur $J$
7236: envisagées dans \hyperref[I.sec9.21.20]{9.21.20} sont satisfaites si $J$ est grand devant
7237: $\Pi'_{\circ}$, stable par $\varphi$ et tel que $\card J\leq \Pi$,
7238: et on achève la démonstration de \hyperref[I.lem9.21.17]{9.21.17} comme en
7239: \hyperref[I.sec9.21.15]{9.21.15}.
7240: 
7241: \begin{enonce*}{Théorème 9.22}\etiquette[9.22]{I.them9.22}
7242: Soit\pageoriginale $p:E\rightarrow B$ un foncteur fibrant, où $B$
7243: est une catégorie essentiellement équivalente à une petite
7244: catégorie, et où pour toute flèche $f:\alpha\rightarrow \beta$ dans
7245: $B$ le foncteur image inverse $f^\ast:E_{\beta}\rightarrow
7246: E_{\alpha}$ est accessible (\hyperref[I.def9.2]{9.2}). Supposons de plus que pour tout
7247: $\alpha\in \ob E$, la catégorie fibre $E_{\alpha}$ admette une
7248: filtration cardinale (\hyperref[I.def9.12]{9.12}) et que tout objet de $E_{\alpha}$ soit
7249: accessible (\hyperref[I.def9.3]{9.3}). Alors chacune des catégories
7250: $F=\Hhom_B(B,E)$ et
7251: $F'=\Hhomcart_B(B,E)$ admet une petite sous-catégorie pleine
7252: génératrice par épimorphismes stricts (\hyperref[I.def7.1]{7.1}), et
7253: chacun de ses objets est accessible.
7254: \end{enonce*}
7255: 
7256: Il est immédiat que l'énoncé ne change pas essentiellement quand on
7257: remplace $B$ par une sous-catégorie pleine $B_{\circ}$ telle que le
7258: foncteur d'inclusion $B_{\circ}\rightarrow B$ soit une équivalence,
7259: et $E$ par $E_{\circ}=E\times_B B_{\circ}$. Cela nous permet de
7260: supposer que $B$ est une petite catégorie.
7261: 
7262: Soit, pour tout $\alpha\in \ob B$, $(\Filt^{\Pi}(E))_{\Pi\geq
7263: \Pi_{\alpha}}$ une filtration cardinale de $E_{\alpha}$. Soit
7264: $c_{\circ}\in 𝒰$ un cardinal satisfaisant les conditions
7265: suivantes:
7266: \begin{equation}
7267: \begin{cases}
7268: & c_{\circ}\geq {\displaystyle{\mathop{\Sup}_{\alpha\in ob
7269: B}}}\Pi_{\alpha}{\quav} c_{\circ}\geq \card %%
7270: \mathrm{F}{}l\, B;\\
7271: & \mbox{pour toute flèche $f:\alpha\rightarrow \beta$ de $B$,
7272: $f^\ast:E_{\beta}\rightarrow E_{\alpha}$ est
7273: $c_{\circ}$-accessible;}\\
7274: & \mbox{pour toute $\alpha\in \ob B$, les objets de
7275: $\Filt^{\Pi_{\alpha}}(E_{\alpha})$ sont $c_\circ$-accessibles.}
7276: \end{cases}\tag{9.22.1}
7277: \end{equation}
7278: Posons
7279: $$
7280: c=2^{c_{\circ}}{\quoi},
7281: $$
7282: de sorte que $c>c_{\circ}$, et pour tout $\alpha\in \ob B$, soit
7283: $$
7284: C_{\alpha}=\Filt^c(E_{\alpha}){\quoi}.
7285: $$
7286: Je\pageoriginale dis que les conditions préliminaires à
7287: \hyperref[I.lem9.21.9]{9.21.9} et
7288: \hyperref[I.lem9.21.17]{9.21.17} sont vérifiées, en faisant
7289: $\Pi=\Pi_{\circ}=c$. C'est clair
7290: pour (9.21.2) et (9.21.3). Pour (9.21.4 a)),
7291: cela résulte de \hyperref[I.cor9.17]{9.17},
7292: et (9.21.4 b)) résulte de la condition \hyperref[I.def9.12]{9.12} c)
7293: pour une filtration
7294: cardinale. Notons d'ailleurs que la condition \hyperref[I.def9.12]{9.12} b) des
7295: filtrations cardinales implique que pour tout $\alpha\in \ob B$, on
7296: a $C^{\Pi}_{\alpha}=C_{\alpha}=\Filt^C(E_{\alpha})$, où
7297: $C^{\Pi}_{\alpha}$ est défini dans (9.21.6). Par suite,
7298: \hyperref[I.them9.22]{9.22} résulte
7299: de \hyperref[I.lem9.21.9]{9.21.9} et de \hyperref[I.lem9.21.17]{9.21.17}, et de façon plus
7300: précise, on a prouvé le
7301: 
7302: \begin{enonce*}{Corollaire 9.23}\etiquette[9.23]{I.cor9.23}
7303: Sous les conditions de \hyperref[I.them9.22]{9.22}, si $c_{\circ}\in 𝒰$ est un
7304: cardinal satisfaisant aux conditions
7305: (9.22.1), et si $c=2^{c_{\circ}}$, alors la sous-catégorie $F^c$ de
7306: $F$ (\resp ${F'}^c$ de $F'$) formée des $X$ tels que $X(\alpha)\in
7307: \ob \Filt^c(E)$ pour tout $\alpha\in \ob B$ est essentiellement
7308: petite et génératrice par épimorphismes stricts; plus précisément,
7309: tout objet $X$ de $F$ (\resp $F'$) est isomorphe à un objet de la
7310: forme $\varinjlim_I X_i$, où les $X_i$ sont dans $F^c$ (\resp dans
7311: ${F'}^c$), et où $I$ est un ensemble ordonné grand devant $c$.
7312: \end{enonce*}
7313: 
7314: Moyennant des hypothèses légèrement plus fortes dans
7315: \hyperref[I.them9.22]{9.22}, on peut
7316: d'ailleurs préciser considérablement \hyperref[I.them9.22]{9.22} et
7317: \hyperref[I.cor9.23]{9.23}:
7318: 
7319: \begin{enonce*}{Corollaire 9.24}\etiquette[9.24]{I.cor9.24}
7320: Sous les conditions de \hyperref[I.them9.22]{9.22}, supposons que chacune des
7321: catégories
7322: $E_{\alpha}$ $(\alpha\in \ob E)$ est stable par petites limites
7323: inductives filtrantes, et que pour tout $\Pi\geq \Pi_{\alpha}$,
7324: $\Filt^{\Pi}(E_{\alpha})$ est stable par limites inductives
7325: filtrantes indexées par des ensembles ordonnés filtrants $I$ tels
7326: que $\card I\leq \Pi$ (ce qui renforce légèrement la condition
7327: \hyperref[I.def9.12]{9.12} {\rm b)} des\pageoriginale filtrations cardinales). Dans le
7328: cas où c'est $F'$ qu'on considère, supposons de plus que pour toute
7329: flèche $f:\alpha\rightarrow \beta$ de $B$, le foncteur
7330: $f^\ast:E_{\beta}\rightarrow E_{\alpha}$ commute aux petites limites
7331: inductives filtrantes. Soit enfin $c_{\circ}\in 𝒰$ un
7332: cardinal satisfaisant aux conditions (9.22.1), et considérons,
7333: pour tout cardinal $\Pi\geq c=2^{c_{\circ}}$, la sous-catégorie
7334: strictement pleine $F^{\Pi}$ (\resp ${F'}^{\Pi}$) de $F$ (\resp de
7335: $F'$) formée des $X$ tels que l'on ait $X(\alpha)\in
7336: \Filt^{\Pi}(E_{\alpha})$ pour tout $\alpha\in \ob B$. Alors les
7337: sous-catégories envisagées définissent une filtration cardinale
7338: (\hyperref[I.def9.12]{9.12}) de $F$ (\resp de $F'$).
7339: \end{enonce*}
7340: 
7341: Il faut vérifier les conditions a), b), c) de \hyperref[I.def9.12]{9.12}. Les
7342: conditions
7343: a) et b) résultent des conditions analogues pour les filtrations
7344: cardinales données des $E_{\alpha}$, et de \hyperref[I.lem9.21.10]{9.21.10}
7345: (\resp\hyperref[I.lem9.21.18]{9.21.18},
7346: compte tenu de la deuxième des conditions (9.22.1)). Reste à prouver
7347: c), dont la première partie résulte aussitôt de
7348: \hyperref[I.lem9.21.9]{9.21.9} (\resp
7349: \hyperref[I.lem9.21.17]{9.21.17}), en faisant $\Pi'_{\circ}=0$, $\Pi_{\circ}=c$. Il reste à
7350: prouver que si on a deux cardinaux $\Pi'\geq \Pi\geq c$, alors pour
7351: tout $X$ dans $F^{\Pi'}$ (\resp ${F'}^{\Pi'}$) on a $X\simeq
7352: \varinjlim_{i\in K}X_i$, avec les $X_i$ dans $F^{\Pi}$ (\resp
7353: ${F'}^{\Pi}$), $K$ grand devant $\Pi$, et $\card K\leq
7354: {\Pi'}^{\Pi}$. Pour ceci, reprenons les démonstrations de
7355: \hyperref[I.lem9.21.9]{9.21.9}
7356: (ii) et de \hyperref[I.lem9.21.17]{9.21.17} (ii). On peut supposer que chaque $I_{\alpha}$
7357: est de cardinal $\leq \Pi'$ en vertu de la condition \hyperref[I.def9.12]{9.12} c) sur la
7358: filtration cardinale de $E_{\alpha}$, donc on aura
7359: $$
7360: \card I\leq ({\Pi'}^{\Pi})^{\card \ob B}\leq
7361: ({\Pi'}^{\Pi})^c={\Pi'}^{\Pi c}={\Pi'}^{\Pi}{\quoi}.
7362: $$
7363: L'ensemble d'indices $K$ utilisé dans 9.21.15 \resp 9.21.21 est
7364: l'ensemble des parties $J$ de $I$ qui sont filtrantes (\ie grandes
7365: devant $\Pi'_{\circ}=0$), stables par $\varphi$ et telles que $\card
7366: J\leq \Pi$. Il est alors immédiat que l'on a
7367: $$
7368: \card K\leq (\card I)^{\Pi}\leq
7369: ({\Pi'}^{\Pi})^{\Pi}={\Pi'}^{\Pi^2}={\Pi'}^{\Pi}{\quoi},
7370: $$
7371: ce\pageoriginale qui achève la démonstration de 9.24.
7372: 
7373: Pour terminer, il convient de donner un énoncé débarrassé des
7374: hypothèses un peu techniques de \hyperref[I.them9.22]{9.22}, en remplaçant
7375: celles-ci par la conjonction, pour les $E_{\alpha}$, des hypothèses
7376: qui interviennent dans \hyperref[I.prop9.11]{9.11} (qui assurent
7377: l'accessibilité des objets des $E$) et dans \hyperref[I.prop9.13]{9.13} (qui
7378: assurent l'existence d'une filtration cardinale dans $E$):
7379: 
7380: \begin{enonce*}{Corollaire 9.25}\etiquette[9.25]{I.cor9.25}
7381: Soit $p:E\rightarrow B$ un foncteur fibrant, où $B$ est une
7382: catégorie équivalente à une petite catégorie, et
7383: où pour toute
7384: flèche $f:\alpha\rightarrow \beta$ de $B$, le foncteur
7385: $f^\ast:E_{\beta}\rightarrow E_{\alpha}$ est accessible (\hyperref[I.def9.2]{9.2}). On
7386: suppose de plus que, pour tout $\alpha\in \ob B$, la catégorie fibre
7387: $E_{\alpha}$ satisfait aux conditions suivantes:
7388: \begin{enumerate}
7389: \renewcommand{\theenumi}{\alph{enumi}}
7390: \renewcommand{\labelenumi}{\theenumi)}
7391: \item $E_{\alpha}$ admet une petite sous-catégorie pleine,
7392: génératrice par épimorphismes stricts (\hyperref[I.def7.1]{7.1}).
7393: 
7394: \item $E_{\alpha}$ est stable par produits fibrés, par noyaux de
7395: doubles flèches, par sommes de deux objets et par conoyaux de
7396: doubles flèches, enfin par petites limites inductives filtrantes
7397: \footnote[1]{Cette dernière condition étant probablement inutile,
7398: \cf 9.13.2.}.
7399: 
7400: \item Le foncteur $\Ker(u,v)$ sur la catégorie des doubles flèches
7401: de $E_{\alpha}$ est accessible (par exemple, commute aux petites
7402: limites inductives filtrantes).
7403: 
7404: \item Tout épimorphisme strict de $E_{\alpha}$ (\hyperref[I.sec10.2]{10.2}) est un
7405: épimorphisme strict universel.
7406: \end{enumerate}
7407: 
7408: Sous ces conditions, chacune des catégories $F=\Hhom_B(B,E)$ et
7409: $F'=\Hhomcart_B$ $(B,E)$ admet une petite sous-catégorie génératrice
7410: par épimorphismes stricts, et tous ses objets sont accessibles
7411: (\hyperref[I.def9.3]{9.3}).
7412: \end{enonce*}
7413: De\pageoriginale plus, \hyperref[I.cor9.24]{9.24} nous donne:
7414: 
7415: \begin{enonce*}{Corollaire 9.26}\etiquette[9.26]{I.cor9.26}
7416: Sous les conditions de $9.25$, et dans le cas où on considère
7417: $F'=\Hhomcart_B(B,E)$, supposons que les foncteurs
7418: $f^\ast:E_{\beta}\rightarrow E_{\alpha}$ commutent aux petites
7419: limites inductives filtrantes. Alors $F$ (\resp $F'$) admet une
7420: filtration cardinale, qu'on peut expliciter par le procédé de
7421: \hyperref[I.cor9.24]{9.24}.
7422: \end{enonce*}
7423: 